% !TeX root = main.tex
\newpage
\section{Vulkan Force Dynamics Source Flow}

\subsection{Purpose}

This section documents the GLSL-side source flow used by the Vulkan
\texttt{ForceDynamicsSimple} family.  It is the Vulkan companion to
\texttt{ForceDynamicsSource.tex}: the Python source document explains the
portable model functions, while this section explains how scenario wrappers
assemble those functions with cell occupancy, lifecycle gates, and optional
features.

The Vulkan wrapper family currently has three documentation targets:
\begin{itemize}
    \item simple particle and function-wall flow;
    \item reservoir piston flow;
    \item streaming reservoir or freestream flow.
\end{itemize}

This update documents the streaming reservoir path used by
\texttt{FreeStream}.  The simple and reservoir-piston sections remain
placeholders until their wrappers are reviewed in the same way.

\begin{figure}[htbp]
\centering
\resizebox{0.98\textwidth}{!}{\begin{tikzpicture}[
    node distance=0.62cm and 1.0cm,
    every node/.style={font=\small},
    stage/.style={draw, rounded corners=2pt, align=center, minimum width=3.65cm, minimum height=0.62cm},
    branch/.style={draw, diamond, aspect=2.1, align=center, inner sep=1.5pt},
    note/.style={draw, dashed, rounded corners=2pt, align=center, minimum width=3.8cm},
    arrow/.style={-{Latex[length=2mm]}, thick}
]

\node[stage] (computeStart) {Compute pass\\\texttt{FreeStreamP.comp}};
\node[branch, below=of computeStart] (sourceGate) {valid active\\mobile source?};
\node[stage, right=of sourceGate] (returnSource) {return};
\node[stage, below=of sourceGate] (init) {Clear diagnostics\\\texttt{totalForce=0}};
\node[stage, below=of init] (dupA) {Initialize duplicate list\\\texttt{MAX\_CELL\_OCCUPANY * 8}};
\node[stage, below=of dupA] (scanA) {Scan source corner cells\\read \texttt{clink}};
\node[branch, below=of scanA] (targetKind) {target kind};
\node[stage, left=of targetKind] (wall) {Boundary marker\\function wall contact};
\node[stage, right=of targetKind] (particle) {Mobile target\\particle contact};
\node[stage, below=of targetKind] (force) {Accumulate contact force\\into source-local total};
\node[stage, below=of force] (velocity) {\texttt{CalcVelocity}\\write alternate velocity};
\node[stage, below=of velocity] (projA) {\texttt{ApplySourceMaximumDepth}\\phase 9401};
\node[stage, below=of projA] (savedContacts) {Saved contact records\\wall and particle geometry};
\node[stage, below=of savedContacts] (scanB) {Replay saved contacts\\register remaining-depth constraints};
\node[stage, below=of scanB] (projB) {\texttt{ApplySourceMaximumDepth}\\phase 9402};
\node[stage, below=of projB] (position) {\texttt{CalcPosition}\\write alternate position};
\node[stage, below=of position] (death) {Cell-domain check\\death bounds update};

\node[note, right=1.1cm of scanB] (perf) {Optimization target:\\one broad-phase scan\\saved contact replay};
\node[note, left=1.1cm of computeStart] (noPiston) {FreeStream excludes\\simple piston include};

\draw[arrow] (computeStart) -- (sourceGate);
\draw[arrow] (sourceGate) -- node[above] {no} (returnSource);
\draw[arrow] (sourceGate) -- node[right] {yes} (init);
\draw[arrow] (init) -- (dupA);
\draw[arrow] (dupA) -- (scanA);
\draw[arrow] (scanA) -- (targetKind);
\draw[arrow] (targetKind) -- node[above] {boundary} (wall);
\draw[arrow] (targetKind) -- node[above] {mobile} (particle);
\draw[arrow] (wall) |- (force);
\draw[arrow] (particle) |- (force);
\draw[arrow] (force) -- (velocity);
\draw[arrow] (velocity) -- (projA);
\draw[arrow] (projA) -- (savedContacts);
\draw[arrow] (savedContacts) -- (scanB);
\draw[arrow] (scanB) -- (projB);
\draw[arrow] (projB) -- (position);
\draw[arrow] (position) -- (death);
\draw[arrow, dashed] (scanB) -- (perf);
\draw[arrow, dashed] (noPiston) -- (computeStart);

\end{tikzpicture}
}
\caption{High-level Vulkan FreeStream compute flow for one source invocation.}
\label{fig:vulkan-force-dynamics-flow}
\end{figure}

\clearpage

\subsection{Simple Wrapper Placeholder}

The simple wrapper is the ordinary particle and function-wall path.  It uses
the shared simple dynamics include, boundary-particle locality helper, and
function-wall helper.  It does not own scheduled inlet release or piston
motion.  This path should be documented after the current FreeStream
performance investigation identifies which common stages are permanent and
which are scenario-specific.

\subsection{Reservoir Piston Placeholder}

The reservoir piston wrapper adds \texttt{ForceDynamicsSimplePiston.glsl} and
defines \texttt{FORCE\_DYNAMICS\_SIMPLE\_PISTON\_AVAILABLE}.  It is the only
simple-wrapper family member that should process a moving analytic piston
before the ordinary target scan.  FreeStream must not enable this path.

\subsection{Streaming Reservoir: FreeStream}

\subsubsection{Files}

The current FreeStream Vulkan path is built from:
\begin{verbatim}
vulkan/sim/FreeStream/FreeStreamP.vert
vulkan/sim/FreeStream/FreeStreamP.comp
vulkan/sim/FreeStream/boundary.glsl
vulkan/sim/FreeStream/params.glsl
vulkan/sim/common/ForceDynamicsSimple.glsl
vulkan/sim/common/ForceDynamicsSimpleBoundaryParticle.glsl
vulkan/sim/common/ForceDynamicsSimpleFunctionWall.glsl
\end{verbatim}

The generated data source is:
\begin{verbatim}
gbase/GenStreaming.py
cfg_gendata/FreeStream.cfg
cfg_gendata/FreeStream_Py.inc
\end{verbatim}

\subsubsection{Streaming Lifecycle}

FreeStream stores timed inlet release state in \texttt{P[SourceID].Data.w}.
The current GLSL lifecycle convention is:
\begin{itemize}
    \item \(\texttt{Data.w}<0\): dead particle;
    \item \(\texttt{Data.w}=0\): active from frame zero;
    \item \(\texttt{Data.w}>0\): pending until
          \(\texttt{ShaderFlags.frameNum} \ge \texttt{Data.w}\).
\end{itemize}

The compute shader applies this gate before force accumulation:
\begin{verbatim}
if (!IsMobileParticleActiveForDynamics(SourceID))
    return;
\end{verbatim}

The vertex shader applies the corresponding lifecycle test before inserting a
particle into the cell-occupancy lists.  Pending and dead particles clear
their corner list to \texttt{npos}, write a suppressed vertex output, and do
not reserve cell slots.

\subsubsection{Graphics Cell-List Stage}

\texttt{FreeStreamP.vert} performs two jobs during the graphics pass:
rendering and broad-phase cell-list construction.  For each non-null active
particle it:
\begin{enumerate}
    \item clears the particle's eight \texttt{CornerList} entries to
          \texttt{npos};
    \item reads the current position buffer selected by
          \texttt{ShaderFlags.positionBuffer};
    \item builds the particle radius AABB;
    \item maps the AABB corners into cell locations;
    \item deduplicates repeated corner cells;
    \item atomically reserves slots in \texttt{L[cell]};
    \item writes the particle ID into \texttt{clink[cell].idx[slot]}.
\end{enumerate}

This makes graphics a required broad-phase stage for compute.  If a particle
does not pass the vertex lifecycle gate, compute will later see no valid
corner cells for that particle.

\subsubsection{Compute Source Stage}

\texttt{FreeStreamP.comp} launches one invocation per particle index.  The
source invocation flow is:
\begin{enumerate}
    \item source zero clears compute counters and returns;
    \item sources beyond \texttt{ShaderFlags.Ptot} return;
    \item non-mobile, pending, or dead sources return;
    \item source diagnostics are cleared;
    \item \texttt{totalForce} is initialized to zero;
    \item a duplicate target array of
          \texttt{MAX\_CELL\_OCCUPANY * 8} entries is initialized;
    \item all source corner cells are scanned for boundary-marker and mobile
          targets;
    \item valid wall and particle contacts accumulate force into
          \texttt{totalForce};
    \item \texttt{CalcVelocity(SourceID, totalForce)} writes the alternate
          velocity buffer;
    \item \texttt{ApplySourceMaximumDepth(SourceID, 9401u)} projects the
          force-driven candidate velocity when needed;
    \item saved wall and particle contact records are replayed to register
          post-resolution remaining-depth constraints;
    \item \texttt{ApplySourceMaximumDepth(SourceID, 9402u)} applies the second
          projection phase;
    \item \texttt{CalcPosition(SourceID)} writes the alternate position buffer
          and applies domain and death-bound rules.
\end{enumerate}

The wrapper does not include piston support.  It includes only the base simple
dynamics, boundary-particle locality helper, and function-wall helper.

\subsubsection{Performance-Relevant Structure}

The current FreeStream performance risk is visible in the wrapper flow rather
than hidden inside the force law.  Each active source pays for:
\begin{itemize}
    \item clearing a duplicate-target array sized by
          \texttt{MAX\_CELL\_OCCUPANY * 8};
    \item scanning up to eight occupied corner cells;
    \item scanning up to \texttt{MAX\_CELL\_OCCUPANY} entries in each cell;
    \item repeating that duplicate-array initialization and cell scan for the
          second maximum-depth phase.
\end{itemize}

Lowering \texttt{MAX\_CELL\_OCCUPANY} improves performance because it reduces
both the scan upper bound and the duplicate-list initialization size.  That
capacity must still remain large enough for the generated initial and runtime
cell occupancy.  The generator should therefore report the smallest safe
capacity and the wrapper should avoid rediscovering the same contact set when
a compact source-owned contact list can carry first-scan results into the
projection phases.

\subsubsection{Intended FreeStream Optimization}

The current wrapper performs two broad-phase scans because the first pass owns
force accumulation and the second pass owns post-velocity remaining-depth
constraint registration.  The two logical phases are still required, but the
second cell-list traversal is not.  The intended FreeStream wrapper shape is:
\begin{verbatim}
scan source corner cells once
deduplicate targets once
save valid particle and wall contact geometry in local contact records
accumulate force while those contacts are discovered
CalcVelocity(SourceID, totalForce)
ApplySourceMaximumDepth(SourceID, 9401u)
loop saved contact records to register/check remaining-depth constraints
ApplySourceMaximumDepth(SourceID, 9402u)
CalcPosition(SourceID)
\end{verbatim}

This changes the second phase from a repeated broad-phase discovery pass into
a compact narrow-phase replay over the actual contacts found for one source.
The saved records are source-local shader temporaries; they are not persistent
particle state and do not change the source-owned dynamics model.

The wrapper should define a source-local contact capacity that represents the
maximum valid simultaneous contacts for one source.  If the saved-contact
array overflows, the shader should report a contact-list capacity error
rather than silently dropping contacts or falling back to a second broad-phase
scan.

\subsubsection{Current Diagnostic Notes}

Shader \texttt{debugPrintfEXT} output is useful for this wrapper, but it must
be treated as diagnostic text rather than a runtime stop condition.  The
Vulkan debug callback should intercept \texttt{VVL-DEBUG-PRINTF}, print only
the message payload, and avoid setting \texttt{Extflg}.

The compute counter \texttt{collOut.numParticles} is incremented after the
active-mobile lifecycle gate when \texttt{DEBUG} is enabled.  This is the
preferred coarse diagnostic for confirming that timed-release particles have
entered compute without using frame-stopping shader printf output.
